\chapter{Theoretical Background}

\section{Navigation in Mobile Robotics}

Navigation is the ability of a robot to act on its knowledge and sensor readings
in order to reach goal positions as efficiently and reliably as
possible~\cite{siegwart2004introduction}.
It can be decomposed into five components:

\begin{enumerate}
  \item \textbf{Perception} -- recording and interpreting raw sensor data such
        as point clouds or images.
  \item \textbf{Localisation} -- estimating the robot pose within a known environment.
  \item \textbf{Mapping} -- constructing a spatial model of the environment
        from sensor observations.
  \item \textbf{Path planning} -- computing a collision-free route from the
        current position to a goal.
  \item \textbf{Motion control} -- generating actuator commands that follow
        the planned path while reacting to unexpected obstacles.
\end{enumerate}

SLAM combines components 2 and 3, estimating the robot pose and the
map at the same time without a priori knowledge.
This thesis focuses exclusively on SLAM, therefore
path planning and motion control are not evaluated. The robot trajectory is
pre-recorded and replayed from a dataset, so no online planning is required.
This also means that the work does not address SPLAM (Simultaneous Planning,
Localisation and Mapping), which links SLAM with active path planning.
The sub-problems relevant to this thesis are described below.

\subsection{Mapping}

Mapping is the problem of building a spatial model $\mathbf{m}$ of the
environment from a sequence of sensor measurements $\mathbf{z}_{1:t}$, under
the assumption that the robot trajectory $\mathbf{x}_{1:t}$, i.e.\ the
sequence of poses, is already known~\cite{thrun2005probabilistic}:

\begin{equation}
  p(\mathbf{m} \mid \mathbf{z}_{1:t},\, \mathbf{x}_{1:t}).
  \label{eq:mapping}
\end{equation}

Even with exact poses the problem is non-trivial: the space of possible maps
is high-dimensional, and repeated sensor noise combined with visually similar
places can lead to incorrect data associations.
Any pose error further deteriorates the result, since a measurement placed at
the wrong location introduces geometrical distortion that spreads across the entire map.
Common map representations used in this thesis are the 2D occupancy grid and
the 3D point cloud, described in Section~\ref{sec:map_representations}.

\subsection{Localisation}

Localisation is the complementary problem to mapping: given a known
map $\mathbf{m}$, the robot determines its current pose $\mathbf{x}_t$ within
that map by fusing sensor measurements $\mathbf{z}_{1:t}$ with
motion commands $\mathbf{u}_{1:t}$~\cite{thrun2005probabilistic}:

\begin{equation}
  p(\mathbf{x}_t \mid \mathbf{m},\, \mathbf{z}_{1:t},\, \mathbf{u}_{1:t}).
  \label{eq:localisation}
\end{equation}

Because pose cannot be read from a sensor directly, the estimate must be
built incrementally: each new measurement updates the previous estimate rather
than recomputing the pose from scratch.
Two operating modes are distinguished depending on prior knowledge~\cite[p.~159]{thrun2005probabilistic}.
In \emph{pose tracking} estimate of starting position is available,
so the search is confined to a small neighbourhood of the current estimate.
This knowledge keeps computational cost low but provides no recovery mechanism if the
robot becomes lost.
In \emph{global localisation} the starting position is completely unknown, so
the belief is initialised uniformly over the entire map and converges toward
the true pose as more observations are collected.
% A harder variant, the \emph{kidnapped robot problem}, arises when the robot is
% moved to an unknown location during operation without being informed, requiring
% the system to detect the inconsistency and re-localise from scratch.
% All algorithms evaluated in this thesis operate in pose tracking mode: they
% process sensor data recorded along a fixed trajectory from a known starting
% position, so global localisation is not required.

\subsection{SLAM}

SLAM assumes that neither the robot poses nor the map are known, and solves
both problems jointly from raw sensor measurements $\mathbf{z}_{1:t}$ and
motion commands $\mathbf{u}_{1:t}$~\cite[p.~245]{thrun2005probabilistic}.
This creates a circular dependency: an accurate map requires accurate poses, and
accurate poses require an accurate map.
The key mechanism that breaks this cycle is loop closure, where recognising a
previously visited place adds a long-range constraint that corrects accumulated
drift across the entire trajectory.

Two problem formulations are common~\cite{thrun2005probabilistic}.
Online SLAM discards past poses and maintains only the current pose
together with the map:
\begin{equation}
  p(\mathbf{x}_t,\, \mathbf{m} \mid \mathbf{z}_{1:t},\, \mathbf{u}_{1:t}).
  \label{eq:slam_online}
\end{equation}
Full SLAM keeps the entire trajectory in the state and optimises the joint
posterior over all poses and the map:
\begin{equation}
  p(\mathbf{x}_{1:t},\, \mathbf{m} \mid \mathbf{z}_{1:t},\, \mathbf{u}_{1:t}).
  \label{eq:slam_full}
\end{equation}
The online formulation is recovered from the full one by integrating out past
poses.
In both cases, the pose is propagated between measurement updates via
$\mathbf{x}_t = \mathbf{x}_{t-1} \oplus \mathbf{u}_t$, where $\oplus$ denotes
composition in $SE(3)$, and each new observation corrects the accumulated drift.
The two dominant solution families, filter-based and graph-based methods, are
described in Sections~\ref{sec:filters} and~\ref{sec:graph_slam} respectively.

% ─────────────────────────────────────────────────────────────────────────────
\section{Sensors}

Sensors provide the data necessary for SLAM implementation. This data is used to
estimates the robot pose and to construct the map.
This thesis uses four sensor types: a 3D LiDAR, a stereo camera, an IMU, and wheel
odometry encoders, each described in the subsections below.

\subsection{LiDAR}

A LiDAR (Light Detection And Ranging) sensor measures distances by emitting laser
pulses and recording the round-trip travel time $\Delta t$.
The distance to the reflecting surface is

\begin{equation}
  d = \frac{c \cdot \Delta t}{2},
  \label{eq:lidar_tof}
\end{equation}

\noindent where $c$ is the speed of light.
A rotating 3D LiDAR fires pulses across multiple vertical channels while spinning
horizontally; each full revolution produces a cloud of tens of thousands of
points described with $(x, y, z)$ coordinates as a
\texttt{sensor\_msgs/PointCloud2} message.
Its main advantages for SLAM are accuracy and independence from ambient lighting.
A relevant limitation is motion distortion: because points within one scan are
acquired at slightly different times while the robot moves, the cloud is
geometrically inconsistent unless corrected.

% A LiDAR (Light Detection And Ranging) sensor measures distances by emitting laser
% pulses and recording the time it takes for each pulse to return after reflecting
% off a surface. This time-of-flight principle gives a direct, metric distance
% measurement for each pulse. A rotating 3D LiDAR fires pulses across multiple
% vertical channels simultaneously while spinning horizontally, producing a dense
% 3D point cloud with each full rotation. A typical scan contains tens of thousands
% of points and is repeated at 10--20\,Hz.
%
% The main advantages of LiDAR for SLAM are its metric accuracy, long range, and
% independence from ambient lighting. A relevant limitation is motion distortion:
% because points in one scan are acquired at slightly different times while the robot
% is moving, the resulting cloud is geometrically inconsistent unless the distortion
% is corrected. LiDAR also provides no colour or texture information, which limits
% appearance-based place recognition.

\subsection{Stereo Camera}

A stereo camera consists of two horizontally separated cameras with a known
baseline $B$ between them.
When the same scene point is visible in both images it appears at slightly
different horizontal positions; this horizontal offset $d$ allows the depth $Z$ to be computed as

\begin{equation}
  Z = \frac{f \cdot B}{d},
  \label{eq:stereo_depth}
\end{equation}

\noindent where $f$ is the focal length in pixels.
Synchronised image pairs are published as \texttt{Image} messages
alongside calibration data in \texttt{CameraInfo}
(both in the \texttt{sensor\_msgs} package).
Cameras provide rich texture information useful for feature matching and
place recognition.
A key limitation is that performance degrades in low-light or
textureless, repeating environments.

% A stereo camera consists of two horizontally separated image sensors with a known
% baseline distance between them. When the same scene point is observed in both
% images, it appears at slightly different horizontal positions (the disparity).
% The depth of the point is proportional to the product of the focal length and
% the baseline, and inversely proportional to the disparity. This gives a dense or
% semi-dense depth image at video frame rate without emitting any signal, making
% stereo cameras passive and unaffected by multi-sensor interference.
%
% Depth uncertainty grows quadratically with distance: at large range, a one-pixel
% disparity error corresponds to a large depth error. Cameras also provide rich
% texture information useful for feature matching and place recognition, but their
% performance degrades in low-light conditions and in textureless environments.

\subsection{360-Degree Camera}

A 360-degree (omnidirectional) camera captures the full horizontal field of
view in a single frame, either through a fisheye lens or by stitching multiple
camera feeds into an equirectangular projection.
Compared to a standard perspective camera, the wider coverage reduces the risk
of losing tracked features when the robot rotates sharply, and provides richer
context for place recognition.

Despite these advantages, 360-degree cameras are not evaluated in this thesis
for two reasons.
First, the visual SLAM algorithms considered — ORB-SLAM3 and RTAB-Map — are
designed around the pinhole camera model; supporting omnidirectional optics
requires either remapping the image to a perspective projection, which
discards the wide-angle benefit, or using specialised fisheye variants that
were outside the scope of this work.
Second, the simulation environment used (Isaac Sim) does not provide a
native omnidirectional camera model, making it impractical to generate
representative training and evaluation data for such a sensor.

\subsection{IMU}

An Inertial Measurement Unit (IMU) combines a three-axis accelerometer and a
three-axis gyroscope.
Both outputs are corrupted by a slowly varying bias $\mathbf{b}$ and Gaussian
noise $\boldsymbol{\eta}$:

\begin{equation}
  \tilde{\mathbf{a}} = \mathbf{a} + \mathbf{b}_a + \boldsymbol{\eta}_a, \qquad
  \tilde{\boldsymbol{\omega}} = \boldsymbol{\omega} + \mathbf{b}_g + \boldsymbol{\eta}_g,
  \label{eq:imu_model}
\end{equation}

\noindent where $\tilde{\mathbf{a}}$ and $\tilde{\boldsymbol{\omega}}$ are the
measured acceleration and angular velocity, and $\mathbf{a}$,
$\boldsymbol{\omega}$ are the true values.
These are published as \texttt{sensor\_msgs/Imu} messages with frequency higher other sensors.
The higher update rate makes the IMU valuable for tracking motion between
slower LiDAR or camera updates.
However, integrating noisy, biased signals causes pose errors to grow without
bound over time, so an IMU must be fused with another sensor for corrections.

% An Inertial Measurement Unit (IMU) combines a three-axis accelerometer and a
% three-axis gyroscope. The accelerometer measures specific force (linear
% acceleration minus gravity) along each body axis; the gyroscope measures angular
% velocity. IMUs operate at high rates, typically 100--400\,Hz, giving dense motion
% information between slower sensor updates.
%
% The main drawback of an IMU is integration drift. Both sensors suffer from
% additive noise and a slowly changing bias. Integrating angular velocity to obtain
% orientation, and doubly integrating acceleration to obtain position, causes errors
% to grow without bound over time. In practice, an IMU is therefore used in
% combination with another sensor such as a LiDAR or camera, which provides
% corrections at a lower rate. This tight coupling can also correct LiDAR motion
% distortion by using the high-rate IMU measurements to interpolate the robot pose
% within a single scan.

\subsection{Wheel Odometry}

Wheel encoders measure the angular velocity of each wheel, from which the
linear speeds $v_L$ and $v_R$ of the left and right wheels are calculated.
For a differential-drive robot the pose increment over a time step $\Delta t$
with orientation $\theta$ is

\begin{equation}
  \Delta x = \frac{v_L + v_R}{2}\,\Delta t\cos\theta, \quad
  \Delta y = \frac{v_L + v_R}{2}\,\Delta t\sin\theta, \quad
  \Delta\theta = \frac{v_R - v_L}{L}\,\Delta t,
  \label{eq:wheel_odom}
\end{equation}

\noindent where $L$ is the distance between the wheels.
Integrating these increments over time provides a pose estimate published as a
\texttt{nav\_msgs/Odometry} message containing position, orientation, and their
covariances.
Wheel odometry is lightweight solution, making it a useful baseline for short-term motion estimation.
However, wheel slip and calibration errors in the wheel
radius and wheel distance cause the estimate to drift over time.

% Wheel odometry estimates the robot's pose increment from encoder counts on the
% drive wheels. For a differential-drive robot such as the TurtleBot3 used in this
% thesis, encoder ticks on the left and right wheels are converted to wheel arc
% lengths, and the relative displacement and heading change are computed from their
% difference. Integrating these increments over time yields a pose estimate in the
% odometry frame.
%
% Wheel odometry is cheap and runs at high frequency, making it a useful short-term
% motion prior. Its accuracy is limited by wheel slip, uneven terrain, and
% systematic errors such as imprecise wheel radius or axle width calibration. These
% errors cause the pose estimate to drift, so wheel odometry alone is not sufficient
% for SLAM but is commonly used as a loose-coupling source alongside more accurate
% sensors.
% For a differential-drive robot such as the TurtleBot3 used in this thesis,
% the pose increment at heading $\theta$ is
%
% \begin{equation}
%   \Delta x = \frac{\Delta s_L + \Delta s_R}{2}\cos\theta, \quad
%   \Delta y = \frac{\Delta s_L + \Delta s_R}{2}\sin\theta, \quad
%   \Delta\theta = \frac{\Delta s_R - \Delta s_L}{L},
%   \label{eq:wheel_odom}
% ─────────────────────────────────────────────────────────────────────────────
\section{Map Representations}
\label{sec:map_representations}

A map is the spatial model of the environment built by a SLAM system.
The choice of representation affects memory consumption, update speed, and how the
map can be used afterwards.
The algorithms evaluated in this thesis produce one of two types of map: a 2D
occupancy grid or a 3D point cloud.

\subsection{Occupancy Grid}

An occupancy grid divides the 2D plane into a regular grid of square cells, each
storing the probability $p(m_i \mid \mathbf{z}_{1:t}, \mathbf{x}_{1:t})$ that
cell $m_i$ is occupied, given all measurements and poses up to time
$t$~\cite{thrun2005probabilistic}.
Rather than storing the probability directly, cells are updated using the log-odds
representation

\begin{equation}
  l_{t,i} = \log \frac{p(m_i \mid \mathbf{z}_{1:t}, 
  \mathbf{x}_{1:t})}{1 - p(m_i \mid \mathbf{z}_{1:t}, \mathbf{x}_{1:t})},
  \label{eq:log_odds}
\end{equation}

% \noindent The log-odds form keeps the update additive and numerically stable.
% At query time the probability is recovered as
% $p = 1 - (1 + e^{l})^{-1}$.

Cells with log-odds above a threshold are marked occupied (obstacles), those below a
second threshold are marked free, and the rest remain unknown.
The main advantage of this representation is that it directly supports path
planning, as each cell already indicates whether the corresponding region is
available for the robot.
Its limitations are that it is restricted to 2 dimensions and that the fixed cell resolution bounds 
the level of detail of the map.

\subsection{3D Point Cloud}

A 3D point cloud is an unordered set of points in a common reference frame,
each defined by its $(x, y, z)$ coordinates.
A global map is built by registering each incoming LiDAR scan into the map
frame using the estimated robot pose and accumulating the resulting points.
Obtained map captures the full three-dimensional geometry of the environment.

Point clouds are well suited for geometric inspection and for evaluating map
quality against a reference scan.
Their drawbacks are high memory consumption in large environments and the
absence of information about free-space: unlike an occupancy grid, a point cloud
records only observed surfaces, not which regions are reachable.

Visual SLAM methods do not accumulate dense range measurements.
Instead, they build sparse point clouds by triangulating matched image
features into 3D landmarks.
The density of such a map depends on the visual texture of the environment,
making it unsuitable for quantitative geometric evaluation.
Visual SLAM maps in this thesis are therefore not evaluated geometrically.

% ─────────────────────────────────────────────────────────────────────────────
\section{Point Cloud Registration}

Point cloud registration is the problem of finding the rigid transformation
that best aligns two overlapping point clouds.
Given a source cloud $\mathcal{P}$ and a target cloud $\mathcal{Q}$, the goal
is to estimate a rotation $\mathbf{R} \in SO(3)$ and translation
$\mathbf{t} \in \mathbb{R}^3$ such that the transformed source
$\{\mathbf{R}\mathbf{p}_i + \mathbf{t}\}$ overlaps with $\mathcal{Q}$ as
closely as possible.
Registration is the core operation in scan-matching-based LiDAR odometry and
SLAM: each new scan is aligned to the previous one or to the accumulated map
to estimate the relative motion of the robot.
The subsections below describe the feature extraction step used to reduce
computation, followed by the four registration algorithms relevant to this
thesis.

\subsection{LiDAR Feature Extraction}

Raw point clouds contain tens of thousands of points per scan, most of which
carry redundant geometric information.
Many algorithms therefore extract a compact set of geometrically distinctive
features before registration.
The local smoothness of a point $i$ is estimated from its neighbours as

\begin{equation}
  c_i = \frac{1}{\lvert \mathcal{S}_i \rvert \cdot \lVert \mathbf{p}_i \rVert}
        \left\lVert \sum_{j \in \mathcal{S}_i} (\mathbf{p}_j - \mathbf{p}_i) \right\rVert,
  \label{eq:smoothness}
\end{equation}

\noindent where $\mathcal{S}_i$ is the set of neighbouring points.
Points with high curvature $c_i$ lie on edges or sharp depth discontinuities;
points with low curvature lie on planar surfaces.
A fixed number of the highest-scoring points are retained as edge features and
the lowest-scoring as planar features.
This selection reduces computation and concentrates matching on the regions
that are most informative for pose estimation.
Feature-based extraction of this form is used by FAST-LIO2 and LIO-SAM,
both evaluated in this thesis.

\subsection{Iterative Closest Point (ICP)}

Iterative Closest Point (ICP) alternates between two steps until convergence.
First, each source point $\mathbf{p}_i$ is matched to its nearest neighbour
$\mathbf{q}_i$ in the target cloud.
Second, the transformation $(\mathbf{R}, \mathbf{t})$ that minimises the
sum of squared point-to-point distances

\begin{equation}
  E(\mathbf{R}, \mathbf{t}) = \sum_i \lVert \mathbf{R}\mathbf{p}_i + \mathbf{t} - \mathbf{q}_i \rVert^2
  \label{eq:icp}
\end{equation}

\noindent is computed in closed form via singular value decomposition.
The process repeats until the change in $E$ falls below a threshold.
ICP is sensitive to initialisation: a poor starting estimate can cause
convergence to a local minimum.
Nearest-neighbour queries are accelerated with a k-d tree, giving
$\mathcal{O}(n \log n)$ cost per iteration.

\subsection{Point-to-Plane ICP}

Point-to-plane ICP replaces the point-to-point error with the distance from
each source point to the tangent plane at its target correspondence,
estimated from the local surface normal $\hat{\mathbf{n}}_i$:

\begin{equation}
  E(\mathbf{R}, \mathbf{t}) =
    \sum_i \bigl( (\mathbf{R}\mathbf{p}_i + \mathbf{t} - \mathbf{q}_i)
                  \cdot \hat{\mathbf{n}}_i \bigr)^2.
  \label{eq:icp_plane}
\end{equation}

\noindent This formulation allows sliding motion along a surface without
incurring a penalty, which makes it more informative on smooth planar
regions and leads to faster convergence than the point-to-point variant.

\subsection{Generalized ICP (GICP)}

Generalized ICP (GICP) extends the point-to-plane idea into a probabilistic
framework by modelling each point as a Gaussian distribution over its local
surface geometry.
The alignment cost becomes the sum of Mahalanobis distances between
corresponding source and target distributions, which subsumes both the
point-to-point and point-to-plane objectives as special cases.
This formulation is more robust to sensor noise and varying point density,
and is the default registration backend of KISS-ICP, evaluated in this
thesis.

\subsection{Normal Distributions Transform (NDT)}

The Normal Distributions Transform (NDT) avoids explicit point-to-point
correspondences entirely.
The target cloud is subdivided into a regular voxel grid, and the points
within each voxel are summarised by their mean $\boldsymbol{\mu}_k$ and
covariance $\boldsymbol{\Sigma}_k$.
Registration maximises the likelihood of the transformed source points under
this mixture of Gaussians:

\begin{equation}
  \text{score}(\mathbf{R}, \mathbf{t}) =
    \sum_i \exp\!\left(
      -\frac{1}{2}(\mathbf{R}\mathbf{p}_i + \mathbf{t} - \boldsymbol{\mu}_k)^\top
      \boldsymbol{\Sigma}_k^{-1}
      (\mathbf{R}\mathbf{p}_i + \mathbf{t} - \boldsymbol{\mu}_k)
    \right),
  \label{eq:ndt}
\end{equation}

\noindent where $k$ is the voxel containing the transformed point.
The fixed voxel structure gives predictable memory usage and fast lookup,
and the method is more robust to non-uniform point density than ICP variants.
NDT is the registration algorithm used by MOLA Odometry, evaluated in this
thesis.

% ─────────────────────────────────────────────────────────────────────────────
\section{Filter-Based State Estimation}
\label{sec:filters}

% TODO: Section intro
%   - Explain that filter-based methods maintain a probability distribution over the robot state
%   - Update it recursively as new sensor data arrives (predict + update cycle)

\subsection{Extended Kalman Filter (EKF)}

% TODO:
%   - State assumptions: Gaussian state distribution, linearised motion and observation models
%   - Describe predict step (propagate mean and covariance via Jacobians) and update step
%   - Note scalability issue: map state grows with number of landmarks → O(n²) cost
%   - Mention use in sensor-fusion odometry (e.g.\ robot\_localization in this thesis)

\subsection{Iterated Extended Kalman Filter (iEKF)}

% TODO:
%   - Explain the limitation of standard EKF: single linearisation at the prior mean
%     can be inaccurate for highly nonlinear systems
%   - Describe the iEKF fix: after the update step, re-linearise around the new state estimate
%     and repeat the update until convergence (typically 2--5 iterations)
%   - Note that this improves accuracy for systems with significant nonlinearity (e.g.\ fast rotation)
%   - Note its use in FAST-LIO2, which employs iEKF on the manifold (SO(3) state)

\subsection{Rao-Blackwellized Particle Filter}

% TODO:
%   - Explain factorisation: each particle carries an independent map; pose is sampled
%   - Describe the SIR update loop: propagate, weight, resample
%   - Note that this is the basis of GMapping (2D LiDAR SLAM) evaluated in this thesis

% ─────────────────────────────────────────────────────────────────────────────
\section{Graph-Based SLAM}
\label{sec:graph_slam}

Most SLAM systems are structured as a two-stage pipeline.
The \emph{front end} processes raw sensor data to extract relative pose
constraints — through scan matching, visual feature tracking, or IMU
integration — and identifies loop closure candidates.
The \emph{back end} takes these constraints as input and optimises the full
trajectory and map for global consistency, typically by solving a nonlinear
least-squares problem over a pose or factor graph.
In contrast to filter-based methods, which process measurements sequentially
and discard past information, graph-based approaches retain the full trajectory
and can enforce global consistency whenever a loop closure is detected.

\subsection{Pose Graphs}

% TODO:
%   - Define nodes (robot poses) and edges (relative pose constraints from odometry or scan-matching)
%   - State the optimisation objective: minimise sum of squared Mahalanobis errors over edges
%   - Mention solvers: g2o, GTSAM, Ceres

\subsection{Factor Graphs}

% TODO:
%   - Generalise pose graphs: nodes can be any variable (poses, landmarks, IMU biases)
%   - Edges are factors (probability distributions over connected variables)
%   - Note that LIO-SAM and MOLA use factor-graph back-ends

\subsection{Loop Closure Detection}

% TODO:
%   - Explain the problem: recognise previously visited places to add long-range constraints
%   - For LiDAR: scan-context, radius-based candidate search + ICP verification
%   - For visual: bag-of-words (DBoW2/DBoW3) — described in detail in visual SLAM section
%   - Note the impact on global consistency: without loop closure, drift accumulates indefinitely

% ─────────────────────────────────────────────────────────────────────────────
\section{IMU Pre-Integration and Deskewing}

% TODO:
%   - Explain LiDAR motion distortion: points in one scan are acquired at different times
%     while the robot moves, causing the cloud to be "smeared"
%   - Describe deskewing: use IMU or estimated velocity to transform all points to a common timestamp
%   - Explain IMU pre-integration: integrate raw IMU measurements between two keyframes
%     into a single relative pose constraint without re-integrating from scratch on each optimisation step
%   - Note the manifold-aware formulation (rotation on SO(3)) used in modern tight-coupling systems

% ─────────────────────────────────────────────────────────────────────────────
\section{Visual SLAM Fundamentals}

% TODO: Section intro
%   - State that visual SLAM uses camera images as the primary sensor
%   - Note the two main paradigms: feature-based (sparse) and direct/semi-dense
%   - This thesis evaluates feature-based stereo systems (ORB-SLAM3, RTAB-Map visual)

\subsection{Stereo Vision and Triangulation}

% TODO:
%   - Recap stereo geometry: epipolar constraint, disparity–depth relationship
%   - State the depth formula: Z = f·B / d (focal length, baseline, disparity)
%   - Note depth uncertainty grows quadratically with distance

\subsection{Feature Detection and Description}

% TODO:
%   - Explain the keypoint pipeline: detect salient points (corners, blobs), compute descriptor,
%     match across frames
%   - GFTT (Good Features To Track / Shi-Tomasi): selects corners by maximising the smaller
%     eigenvalue of the local gradient matrix; computationally cheap, produces stable 2D tracks
%     across consecutive frames — used as the keypoint detector in RTAB-Map's visual front-end
%   - BRIEF (Binary Robust Independent Elementary Features): encodes a patch as a binary string
%     by comparing intensities at randomly sampled point pairs; Hamming-distance matching is
%     very fast; not rotation- or scale-invariant on its own
%   - ORB (Oriented FAST + Rotated BRIEF): adds orientation estimation to FAST keypoints
%     and steers BRIEF accordingly, giving rotation invariance with minimal extra cost;
%     also applies a scale pyramid for scale invariance; used by ORB-SLAM3
%   - Mention matching strategy: brute-force Hamming distance or FLANN with ratio test

\subsection{Visual Odometry}

% TODO:
%   - Describe the VO pipeline: detect features → match across frames → estimate relative pose via PnP or essential-matrix decomposition
%   - Explain bundle adjustment (local BA) as the front-end optimisation
%   - Note accumulation of drift without loop closure

\subsection{Bag-of-Words Place Recognition}

% TODO:
%   - Explain BoW: cluster descriptor space into a visual vocabulary (k-d tree / k-means)
%   - Each image is represented as a sparse histogram over visual words
%   - Loop closure candidate: high cosine similarity between current and past BoW vectors
%   - Geometric verification with RANSAC before adding the loop-closure edge

% ─────────────────────────────────────────────────────────────────────────────
\section{Evaluation Metrics}
\label{sec:eval_metrics}

Two complementary groups of metrics are used in this thesis.
Trajectory metrics quantify how closely the estimated robot path matches the
ground truth poses recorded by the simulator.
Map metrics quantify the geometric fidelity of the 3D point cloud produced
by each algorithm relative to the ground-truth map extracted from the
simulator.
The ground truth is a time-stamped sequence of $SE(3)$ poses, one per
sensor frame, with sub-millimetre positional accuracy.

\subsection{Umeyama Alignment}

An estimated trajectory may be displaced or rotated relative to the ground
truth because SLAM initialises its coordinate frame from the first sensor
measurement rather than from an absolute reference.
Before computing error statistics, the two trajectories are therefore
brought into the same frame using the method of Umeyama~\cite{umeyama1991least},
which finds the similarity transformation parameters — rotation $\mathbf{R}$,
translation $\mathbf{t}$, and scale $c$ — that minimise the mean squared error

\begin{equation}
  \varepsilon^2(\mathbf{R}, \mathbf{t}, c) =
    \frac{1}{n}\sum_{i=1}^{n}
    \lVert \mathbf{q}_i - (c\mathbf{R}\mathbf{p}_i + \mathbf{t}) \rVert^2,
  \label{eq:umeyama_cost}
\end{equation}

\noindent where $\mathbf{p}_i \in \mathbb{R}^3$ are the estimated positions and
$\mathbf{q}_i \in \mathbb{R}^3$ are the corresponding ground-truth positions.
The optimal rotation is obtained in closed form via singular value decomposition
of the cross-covariance matrix
$\boldsymbol{\Sigma}_{qp} = \frac{1}{n}\sum_i
(\mathbf{q}_i - \bar{\mathbf{q}})(\mathbf{p}_i - \bar{\mathbf{p}})^\top$,
and the optimal scale follows as $c^* = \frac{1}{\sigma_p^2}\operatorname{tr}(\mathbf{D}\mathbf{S})$,
where $\sigma_p^2$ is the variance of the estimated positions,
$\mathbf{D}$ contains the singular values of $\boldsymbol{\Sigma}_{qp}$,
and $\mathbf{S}$ is a sign matrix that ensures a proper rotation.
For metric sensor modalities (LiDAR, stereo), the recovered scale is close to
unity and the alignment reduces to $SE(3)$.
For monocular visual methods, which recover geometry only up to an unknown
scale, the scale $c$ is fully optimised, yielding a $SIM(3)$ alignment.

\subsection{Absolute Trajectory Error (ATE)}

The Absolute Trajectory Error measures the global geometric consistency
of the estimated path~\cite{sturm2012benchmark}.
After Umeyama alignment, the translational error at each time step $i$ is
the Euclidean distance between the aligned estimated position
$\hat{\mathbf{p}}_i$ and the ground truth position $\mathbf{p}_i^*$.
The RMSE over all $n$ poses is

\begin{equation}
  \text{ATE} = \sqrt{\frac{1}{n}\sum_{i=1}^{n}
               \lVert \hat{\mathbf{p}}_i - \mathbf{p}_i^* \rVert^2}.
  \label{eq:ate}
\end{equation}

A large ATE indicates that the algorithm accumulates drift over the full
trajectory, or that loop closure is absent or ineffective.
In addition to the RMSE, the mean and maximum translational errors are
reported to characterise the error distribution.

\subsection{Relative Pose Error (RPE)}

The Relative Pose Error measures local odometric accuracy independently of
the global alignment~\cite{sturm2012benchmark}.
For each consecutive pose pair $(i,\, i{+}1)$ the relative motion is
computed both in the estimate and in the ground truth, and the translational
part of the residual between the two relative poses is taken as the local
error.
The RMSE over all such pairs is

\begin{equation}
  \text{RPE} = \sqrt{\frac{1}{n-1}\sum_{i=1}^{n-1}
               \lVert \mathbf{e}_{i}^{\,\text{trans}} \rVert^2},
  \label{eq:rpe}
\end{equation}

\noindent where $\mathbf{e}_{i}^{\,\text{trans}}$ is the translational
component of the residual between the relative pose in the estimate and the
relative pose in the ground truth over the same one-step window.
Because RPE is computed over short, fixed-length windows it captures
short-term drift that ATE may hide through loop-closure corrections.

\subsection{Chamfer Distance}

The Chamfer Distance (CD) is a symmetric measure of geometric fidelity
between two point clouds~\cite{fan2017point}.
Given an estimated cloud $\mathcal{A}$ and a reference cloud $\mathcal{B}$,
it averages the nearest-neighbour distance in both directions:

\begin{equation}
  \text{CD}(\mathcal{A},\mathcal{B}) =
    \frac{1}{|\mathcal{A}|}\sum_{\mathbf{a}\in\mathcal{A}}
      \min_{\mathbf{b}\in\mathcal{B}}\lVert\mathbf{a}-\mathbf{b}\rVert
    +
    \frac{1}{|\mathcal{B}|}\sum_{\mathbf{b}\in\mathcal{B}}
      \min_{\mathbf{a}\in\mathcal{A}}\lVert\mathbf{b}-\mathbf{a}\rVert.
  \label{eq:chamfer}
\end{equation}

The first term penalises spurious points in the estimate that have no
counterpart in the reference; the second term penalises incomplete regions
of the map.
CD is reported in metres and reflects the overall geometric accuracy of the
reconstructed map, but is sensitive to outlier points, which can inflate the
score without reflecting a systematic reconstruction failure.

\subsection{F-score}

The F-score at a distance threshold $d$ combines precision and recall into a
single bounded quality measure~\cite{knapitsch2017tanks}.
A point $\mathbf{a} \in \mathcal{A}$ is considered correct if its nearest
neighbour in $\mathcal{B}$ is within $d$, and vice versa:

\begin{align}
  P(d) &= \frac{1}{|\mathcal{A}|}
          \bigl|\{\mathbf{a}\in\mathcal{A} :
          \min_{\mathbf{b}}\lVert\mathbf{a}-\mathbf{b}\rVert < d\}\bigr|, \\[4pt]
  R(d) &= \frac{1}{|\mathcal{B}|}
          \bigl|\{\mathbf{b}\in\mathcal{B} :
          \min_{\mathbf{a}}\lVert\mathbf{b}-\mathbf{a}\rVert < d\}\bigr|, \\[4pt]
  F(d) &= \frac{2\,P(d)\,R(d)}{P(d)+R(d)}.
  \label{eq:fscore}
\end{align}

Three thresholds are used in this thesis: $d = 5$\,cm (strict),
$d = 10$\,cm (moderate), and $d = 20$\,cm (lenient).
F@5\,cm penalises any gross geometric errors, while F@20\,cm tolerates
small misalignments and rewards completeness.
Unlike the Chamfer Distance, the F-score is bounded to $[0, 1]$ and is not
dominated by a small number of distant outliers.
